\chapter{Identification before estimation}\label{ch:identification}
Identification determines whether a requested answer is a function of the
available distribution and stated assumptions. Estimation approximates that
function from a finite sample. Computation evaluates an estimator. The three
questions fail in different ways: an observational equivalence cannot be
repaired by a larger cluster, a noisy identified effect may benefit from more
independent observations, and a slow valid estimator may benefit from better
software. A product must diagnose which problem it is facing.
\section{The estimand is a product contract}
For an eligible household let $A\in\{0,1\}$ denote assigned treatment, $X$
information available before assignment, and $Y(a)$ the horizon outcome under
assignment $a$. The average treatment effect is
\begin{equation}
 \theta=\E[Y(1)-Y(0)],\qquad
 \tau(x)=\E[Y(1)-Y(0)\mid X=x].
 \label{eq:ate}
\end{equation}
The expectation refers to a declared population. An effect among existing
high-balance customers is not an effect among all households. Changing
eligibility, outcome horizon or economic valuation creates a new
estimand. For offers, the intention-to-treat effect concerns assignment;
the effect of receipt requires a separate compliance design.

Potential outcomes make the contrast explicit. Structural causal models
describe how interventions replace mechanisms. A single-world intervention
graph links the two by splitting a treatment node into its fixed intervention
value and its random component \citep{hernan2020,peters2017,swig2013}.
The product should store the estimand and, where used, a graph with
substantive variable definitions. Neither notation substitutes for a plausible
assignment mechanism.

Under consistency, conditional exchangeability
$\{Y(0),Y(1)\}\perp A\mid X$, and positivity
$0<e(X)=\Prob(A=1\mid X)<1$ on the target population,
\begin{equation}
 \theta=\E[m_1(X)-m_0(X)],\qquad m_a(x)=\E[Y\mid A=a,X=x].
 \label{eq:adjust}
\end{equation}
Consistency includes the absence of consequential undisclosed treatment
versions. An offer is inadequately defined if its price, delivery and duration
are unspecified. Exchangeability concerns actual decision information:
a relationship manager's unrecorded knowledge can invalidate it even if the
warehouse contains thousands of predictive fields. Positivity concerns actions
within histories; large row counts do not create support for actions never taken.

\section{Experiments and observational studies}
Randomize at household level when accounts share money, or at branch/region
level when capacity or communications spill across customers. Preserve actual
allocation probabilities and distinguish assignment, delivery and take-up.
Experiments also need a prespecified horizon, contamination diagnostics and
attrition analysis.

Adaptive experiments require their time-varying assignment probabilities and
inference valid for the adaptation and stopping rule. Repeatedly examining
ordinary fixed-sample confidence intervals does not make them valid under
optional stopping. In a cluster experiment, the number of randomized
clusters constrains inference even with millions of transactions per cluster.
Randomization inference uses the actual allocation mechanism, not an
analyst's preferred regression model.

Observational adjustment is appropriate when assignment information and
overlap are defensible. Diagnostics include balance, propensity support,
outcome sensitivity and negative controls with justified causal roles.
They can discover defects but do not prove that all confounders were observed.
The uncertainty object should distinguish sampling uncertainty from
assumptions about omitted confounding \citep{ovb2022}.

\section{Why orthogonal estimation is useful}
Flexible prediction helps estimate nuisance functions $m_a$ and $e$, but
direct substitution can introduce regularization bias. An augmented
inverse-probability score is
\begin{equation}
 \phi(O;m,e)=m_1(X)-m_0(X)
 +\frac{A\{Y-m_1(X)\}}{e(X)}
 -\frac{(1-A)\{Y-m_0(X)\}}{1-e(X)}.
 \label{eq:aipw}
\end{equation}
Estimate nuisances outside the evaluation fold and average held-out scores.
Cross-fitting limits direct dependence between nuisance overfitting and
evaluation residuals. Orthogonality reduces first-order sensitivity to
nuisance error \citep{dml2018,wager2025}.

\begin{proposition}[Conditional bias of the AIPW score]
With true nuisances $e,m_0,m_1$, fixed candidate nuisances
$\widetilde e,\widetilde m_0,\widetilde m_1$ and denominators bounded away
from zero,
\begin{align}
&\E[\phi(O;\widetilde m,\widetilde e)\mid X]-(m_1-m_0)(X)\nonumber\\
&\quad=(\widetilde e-e)
 \left\{\frac{\widetilde m_1-m_1}{\widetilde e}
       +\frac{\widetilde m_0-m_0}{1-\widetilde e}\right\}(X).
 \label{eq:drbias}
\end{align}
\end{proposition}
\begin{proof}
Conditioning on $X$ replaces the treated residual term by
$e(m_1-\widetilde m_1)/\widetilde e$ and the untreated residual term by
$-(1-e)(m_0-\widetilde m_0)/(1-\widetilde e)$. Collecting coefficients of
each outcome error gives the expression.
\end{proof}

The bias vanishes if either nuisance component is correct; otherwise it
depends on products of errors. With overlap bounds, Cauchy--Schwarz bounds
these terms by products of nuisance-error norms. Under the remaining
independent-sampling and moment conditions, products
$o_p(n^{-1/2})$ support asymptotically linear inference. This is not a
result for arbitrarily dependent transaction rows.

For independent observations, the empirical held-out score variance divided
by $n$ estimates the variance of the mean under suitable conditions.
Household trajectories and common shocks require appropriate score
construction and covariance aggregation. The interface therefore requires
both an estimator and an inference specification.

\section{A campaign example with a known answer}
The synthetic construction sets $X\sim\mathrm{Bernoulli}(1/2)$,
$e(X)=0.2+0.6X$, and
\[
 Y=100+80X+(10-4X)A+\varepsilon,\qquad
 \varepsilon\sim N(0,20^2).
\]
The true effect is $8$. In the population,
$\Prob(X=1\mid A=1)=0.8$ and $\Prob(X=1\mid A=0)=0.2$.
The observed difference is $80(0.8-0.2)+10-4(0.8)=54.8$:
assignment targets customers with larger untreated outcomes.

In 40,000 independent synthetic households, five-fold saturated nuisance
fits give AIPW $\NumStaticATE$ with standard error $\NumStaticSE$,
compared with the raw difference $\NumStaticNaive$.
The 95\% interval is approximately $[7.974,8.969]$.
This checks a computation against a known process. One seed and an interval
containing the truth do not demonstrate coverage.

Heterogeneity matters economically. Effects are $10$ and $6$ in the two
groups. At an additional cost of $7$ per offer, treating everyone has true
incremental value $1$ per eligible customer; treating only $X=0$ has
value $1.5$. The code evaluates these fixed policies with paired AIPW scores
and simultaneous intervals over three declared comparisons.
This demonstrates policy comparison without optimizing and evaluating a
learned policy on the same outcomes.
The targeted-versus-all estimate is $0.336$ with a simultaneous interval
approximately $[-0.096,0.769]$. Although the generating model makes the
targeted policy better by $0.5$, this sample does not establish superiority
at the declared simultaneous confidence level. Both policies have positive
lower bounds relative to treating nobody. The example thus distinguishes
positive value from evidence that a more selective policy is better.

\section{Heterogeneity, instruments and thresholds}
Meta-learners and orthogonal learners estimate conditional average effects
\citep{kunzel2019,dml2018}. Compare subgroup calibration, policy value and
uncertainty, not only factual prediction. Suppress unsupported segments.
A customer's CATE is an average among comparable customers, not an
observable personal causal effect.

When assignment has unmeasured determinants, a valid instrument can offer
a different route \citep{applied2024}. With binary randomized encouragement $Z$, exclusion,
relevance, independence and monotone take-up identify the complier effect
through the Wald ratio
\begin{equation}
 \theta_{\mathrm{LATE}}=
 \frac{\E[Y\mid Z=1]-\E[Y\mid Z=0]}
      {\E[A\mid Z=1]-\E[A\mid Z=0]}.
 \label{eq:wald}
\end{equation}
Monotonicity and exclusion are not consequences of a strong first stage.
Weak first stages need suitable inference. Competitor rates may directly
affect switching and therefore fail exclusion. A price instrument valid in
one market can fail in another because information and policy differ.

Regression discontinuity uses a treatment rule at a threshold $c$ of a
running variable $R$. With continuity of potential-outcome means, the
outcome jump identifies a local effect in a sharp design. In a fuzzy design
divide the outcome jump by the treatment-probability jump, with appropriate
compliance restrictions. Retain the score and rule version used at the time;
inspect manipulation, competing thresholds and bandwidth sensitivity.
The local effect does not by itself value a policy far from the threshold.

\section{What a result is allowed to say}
Every result identifies its design, target, comparison, assumptions and
uncertainty. Its scientific state may be \emph{identified under declared
assumptions}, \emph{partially identified}, \emph{unsupported by observed
action overlap}, or \emph{not identified under the supplied model}.
Numerical convergence is a separate field. A consuming application can then
abstain or request a better experiment without mistaking prediction for
identification.

\section{Why the augmented score works}
The AIPW score has a useful interpretation for a banker. It first predicts
each customer's outcome under each action, then corrects the prediction
using the observed residual among customers who actually received that
action. A rare action receives a larger correction because fewer observed
customers represent its potential outcome. The estimator does not create
counterfactual observations; it uses exchangeability to give those residuals
their intervention interpretation.

Equation~\eqref{eq:drbias} makes that forecast-and-correction account
precise. Conditioning on training information treats fitted functions as
fixed for an independent evaluation observation; averaging the conditional
remainder gives the population bias. The relevant protection concerns
the product of nuisance errors, with the overlap and inference conditions
already stated. It does not assert that two inaccurate learners
necessarily cancel.

For bank data, overlap can vary sharply by product, jurisdiction and
eligibility rule. A risk department may never approve a limit increase for
a certain population. No flexible outcome learner makes the intervention
identified there from those records. Trimming changes the population;
capping weights changes the estimator's bias; replacing a deterministic
policy by a feasible stochastic policy changes the estimand. Each can be
reasonable, but the report must describe the actual choice and recompute
the economic target accordingly.

\section{Residualization can change the target}
Partially linear DML is attractive because it combines flexible nuisance
prediction with a low-dimensional coefficient. Its canonical model is
\[
 Y=\theta_0 A+g_0(X)+U,\quad \E[U\mid A,X]=0,
 \qquad A=e_0(X)+V,\quad \E[V\mid X]=0.
\]
Let $\ell_0(X)=\E[Y\mid X]$. Regress the residualized outcome
$Y-\ell_0(X)$ on $A-e_0(X)$, or solve the orthogonal moment
\begin{equation}
 \E[(A-e_0(X))\{Y-\ell_0(X)-\theta(A-e_0(X))\}]=0.
 \label{eq:plr-score}
\end{equation}
Cross-fitting allows flexible estimates of $e_0$ and $\ell_0$ while
retaining a tractable target under the model and rate assumptions.

Now suppose the true effect is $\tau(X)$ rather than a constant. For a
binary action with conditional mean $m_0(X)+A\tau(X)$, the population
solution of the same moment is
\begin{equation}
 \theta_{\mathrm{PLR}}=
 \frac{\E[e_0(X)\{1-e_0(X)\}\tau(X)]}
      {\E[e_0(X)\{1-e_0(X)\}]}.
 \label{eq:plr-weighted}
\end{equation}
To obtain it, condition the numerator
$\E[(A-e_0)(Y-\ell_0)]$ on $X$; it equals the conditional
treatment variance times $\tau(X)$. The denominator averages that same
variance. The resulting coefficient emphasizes customers with treatment
probabilities near one half. It need not equal the eligible population's
ATE or an effect among treated customers.

This distinction is operational. A bank could compare an AIPW ATE to a
partially linear DML coefficient, see different values and incorrectly
attribute the difference to software quality. The estimates may target
different weighted populations. A challenger comparison must either align
the target or state the different question. The adapter registry therefore
needs an estimand declaration, including aggregation weights, rather than
only a label such as ``DML supported.''

\section{Selecting covariates requires temporal and causal reasoning}
A variable measured before an action can still be inappropriate to adjust
for if it is a collider or an instrument that aggravates variance without
addressing confounding. A variable measured after an action can be a mediator,
a common effect or a confounder of a later decision. The temporal clock is
necessary, but it does not decide every adjustment question.

For example, an offer $A_0$ changes engagement $L_1$, which affects both
future assistance $A_1$ and final recoveries $Y$. Adjusting for $L_1$
in a single regression intended to estimate the total $A_0$ effect can
block part of that effect. Failing to adjust for it when evaluating $A_1$
can leave confounding. Chapter~\ref{ch:longitudinal} resolves this through
the sequential law. It cannot be resolved by a global ``pre-treatment
features'' list shared by all decision dates.

The same issue appears in customer service data. Complaints can result from
both an offer and latent dissatisfaction. Conditioning on complaint receipt
may select a subgroup whose latent dissatisfaction differs across offers.
An analysis of customers with a recorded complaint then targets neither
the original cohort nor an automatically valid mediator effect. A mediation
study needs its own intervention definition and identifying assumptions;
the total assignment effect remains useful even when the mechanism is
unresolved \citep{hernan2020,swig2013}.

The proposed data organizer can suggest which columns describe complaints,
assignments and account outcomes. It can identify timestamps and draw
candidate graphs. It cannot infer that a complaint is an acceptable
adjustment variable from its business name. The study must retain the
reason it is included, excluded or treated as a time-varying variable.

\section{Measurement, missingness and transport are part of identification}
Suppose true treatment is an assigned offer but the source records only an
opened message. The measurement is then affected by attention, device use
and customer preference. Relabeling message opening as treatment changes
the intervention. Recovering a hidden treatment from surrogates requires
additional measurement and proxy restrictions; recent work does not justify
universal recovery of missing action logs \citep{hidden2026}.

Likewise, a missing competitor balance is not a zero balance. If observed
bank balance equals liquid wealth times bank share, $B=W\omega$, the
same $B$ is compatible with many $(W,\omega)$ pairs. Neither precise
measurement of $B$ nor a flexible causal learner identifies the outside
wallet without additional information. Data-recovery draws can propagate
uncertainty under a supplied missingness model, but they do not establish
that model. The corrupted-data literature makes clear that different
measurement and observation operators require different identification
conditions \citep{corrupt2022}.

For a source population $S=1$ and a target population $S=0$, transporting
an effect requires more than matching feature names. If conditional causal
effects are stable and every target $X$ has source support, the target
ATE is $\E[\tau(X)\mid S=0]$. This is a useful standardization formula.
It does not follow merely because a classifier has difficulty distinguishing
the populations. Product terms, laws, delivery systems, hidden effect
modifiers and competitive responses may change the response mechanism.
The target population and stability assumptions must remain attached to
the transported estimate.

\section{From an effect to a financial comparison}
Let $C_k(a)$ be dated cash or economic contribution components with signs
defined from the group's perspective, and let $d_k$ be fixed discount and
currency-conversion factors. Define $Y(a)=\sum_k d_k C_k(a)$.
Linearity gives
\[
 \E[Y(1)-Y(0)]=\sum_k d_k\E[C_k(1)-C_k(0)].
\]
However, uncertainty is joint. If $\widehat{\boldsymbol\theta}$ estimates
the component effects with covariance $\widehat\Sigma$, the variance
of the financial contrast is $\boldsymbol d'\widehat\Sigma\boldsymbol d$.
Adding component confidence intervals or treating correlated components as
independent does not reproduce this variance. Fitting all components on
the same folds and retaining customer or cluster score vectors allows
the covariance to be reconstructed.

This formula conditions on fixed financial factors. An uncertain transfer
price, loss severity or capital charge adds parameter and scenario uncertainty.
If the policy changes aggregate market rates, the factors can themselves
be intervention-dependent. Then a common fixed conversion is an approximation
whose adequacy needs an economic argument. Deposit principal is a funding
quantity; its full value is not an immediate profit. Similarly, a loan
advance is an outflow paired with a claim, not an expense that should be
counted again when expected principal losses are recorded.

Scientific uncertainty and economic sensitivity should be displayed separately
and, where a defensible joint model exists, combined. A narrow sampling interval
under one accounting convention does not settle a decision that reverses
under plausible funding-cost or customer-response assumptions.
